% How the application was designed, audited, and measured (report plan): threat
% model, ASVS Level 2 audits, tools, both teammates attacking the whole
% application, the detection study, and how test data was selected.
\section{Methodology}

\subsection{Threat model}
Before writing code, we drew a data-flow diagram with its trust boundaries and applied STRIDE to it with four attacker profiles: P1, an outside attacker with no account; P2, a malicious employee; P3, a compromised manager or administrator; and P4, rogue platform staff with server or database access. Each threat has a stable identifier, a protection, and the test that checks it; audit findings name the threat they fall under. \pending{figure: the data-flow diagram.}

\subsection{Audits}
We will audit against OWASP ASVS Level 2 \cite{asvs5}, applicable chapters only, on local test copies, never on systems we do not own. Both of us will attack the whole application together in every audit rather than splitting it by who wrote what, because people miss flaws in their own code. The planned audits are: the backend, the frontend, an internal-user audit (P2 and P3), and an external-attacker audit (P1). The tools are OWASP ZAP, sqlmap, Burp Suite Community, and the static analysis and dependency checks that run on every pull request. Every finding will be recorded as \emph{tool, vulnerability, code fix, regression test}, with its attacker profile, component, and severity. \pending{switch to the past tense as each audit is done; add its dates and the number of ASVS requirements checked.}

\subsection{Detection study}
The planned detector normalizes each text value the API receives (URL and HTML decoding, Unicode NFKC, lowercasing) and matches it against regular expressions for five attack types: SQL injection, cross-site scripting, path traversal, command injection, and format string or template injection. It only records matches; it never blocks a request. \pending{switch to the past tense once the detector is built.}

\textbf{Data selection.} Attack examples will come from PayloadsAllTheThings and SecLists, pinned to one commit each. Harmless examples will be tricky text written by the team before tuning begins (for example \emph{O'Brien} and ``please select a new shift''), the demo data's text, and a public surname list. Duplicates will be removed after normalization, near-duplicates kept in the same set, and large types sampled down to 1,000 entries with a fixed seed. A field set of requests sent by the audit tools will check the results on traffic the patterns were not built from. \pending{final counts per type, and the excluded files with reasons.}

\textbf{Scoring.} Each type will be split 60\% for tuning and 40\% for testing, stratified by type and source, and the split committed with checksums before tuning starts. The tagged detector will be scored on the test set once. We will report precision, recall, and $F_1$ per type with Wilson 95\% confidence intervals, a confusion matrix, and matching time per value, including crafted worst-case input. \pending{switch data selection and scoring to the past tense once the split is committed and the test set is scored.}

\subsection{Replication}
Datasets, scripts, results, and this report are published in the \texttt{workforce-ops-research} repository \pending{repository URL and the commit the results come from}.
